%% 
%% Copyright 2007-2024 Elsevier Ltd
%% 
%% This file is part of the 'Elsarticle Bundle'.
%% ---------------------------------------------
%% 
%% It may be distributed under the conditions of the LaTeX Project Public
%% License, either version 1.3 of this license or (at your option) any
%% later version.  The latest version of this license is in
%%    http://www.latex-project.org/lppl.txt
%% and version 1.3 or later is part of all distributions of LaTeX
%% version 1999/12/01 or later.
%% 
%% The list of all files belonging to the 'Elsarticle Bundle' is
%% given in the file `manifest.txt'.
%% 
%% Template article for Elsevier's document class `elsarticle'
%% with numbered style bibliographic references
%% SP 2008/03/01
%% $Id: elsarticle-template-num.tex 249 2024-04-06 10:51:24Z rishi $
%%
\documentclass[preprint,12pt]{elsarticle}

%% Use the option review to obtain double line spacing
%% \documentclass[authoryear,preprint,review,12pt]{elsarticle}

%% Use the options 1p,twocolumn; 3p; 3p,twocolumn; 5p; or 5p,twocolumn
%% for a journal layout:
%% \documentclass[final,1p,times]{elsarticle}
%% \documentclass[final,1p,times,twocolumn]{elsarticle}
%% \documentclass[final,3p,times]{elsarticle}
%% \documentclass[final,3p,times,twocolumn]{elsarticle}
%% \documentclass[final,5p,times]{elsarticle}
%% \documentclass[final,5p,times,twocolumn]{elsarticle}

%% For including figures, graphicx.sty has been loaded in
%% elsarticle.cls. If you prefer to use the old commands
%% please give \usepackage{epsfig}

%% The amssymb package provides various useful mathematical symbols
\usepackage{amssymb}
%% The amsmath package provides various useful equation environments.
\usepackage{amsmath}
%% The amsthm package provides extended theorem environments
%% \usepackage{amsthm}

%% The lineno packages adds line numbers. Start line numbering with
%% \begin{linenumbers}, end it with \end{linenumbers}. Or switch it on
%% for the whole article with \linenumbers.
%% \usepackage{lineno}

\journal{Artificial Intelligence in Agriculture}

\begin{document}

\begin{frontmatter}

%% Title, authors and addresses

%% use the tnoteref command within \title for footnotes;
%% use the tnotetext command for theassociated footnote;
%% use the fnref command within \author or \affiliation for footnotes;
%% use the fntext command for theassociated footnote;
%% use the corref command within \author for corresponding author footnotes;
%% use the cortext command for theassociated footnote;
%% use the ead command for the email address,
%% and the form \ead[url] for the home page:
%% \title{Title\tnoteref{label1}}
%% \tnotetext[label1]{}
%% \author{Name\corref{cor1}\fnref{label2}}
%% \ead{email address}
%% \ead[url]{home page}
%% \fntext[label2]{}
%% \cortext[cor1]{}
%% \affiliation{organization={},
%%             addressline={},
%%             city={},
%%             postcode={},
%%             state={},
%%             country={}}
%% \fntext[label3]{}

\title{An Explainable AI-Based Decision Support Framework for Agricultural Skill Development Programs}

%% use optional labels to link authors explicitly to addresses:
%% \author[label1,label2]{}
%% \affiliation[label1]{organization={},
%%             addressline={},
%%             city={},
%%             postcode={},
%%             state={},
%%             country={}}
%%
%% \affiliation[label2]{organization={},
%%             addressline={},
%%             city={},
%%             postcode={},
%%             state={},
%%             country={}}

\author{Kitti Chiewchan} %% Author name

%% Author affiliation
\affiliation{organization={Computer Education, Faculty of Education, Bansomdejchaopraya Rajabhat University},%Department and Organization
            addressline={1061 Soi Issarapab 15, Issarapab Road, Hiranruchi Subdistrict, Thonburi District}, 
            city={Bangkok},
            postcode={10600}, 
            state={Bangkok},
            country={Thailand}}

%% Abstract
\begin{abstract}
%% Text of abstract
The rapid advancement of artificial intelligence (AI) has transformed agricultural systems from input-driven production models into data-driven decision ecosystems. While smart agriculture technologies have primarily focused on optimizing physical processes such as crop monitoring, irrigation, and yield prediction, comparatively less attention has been devoted to the development of AI-driven systems that support human capital formation — a critical determinant of long-term agricultural productivity and resilience.

Agricultural skill development programs represent one of the most widely implemented policy instruments for improving farmer performance. However, their effectiveness remains limited by a persistent structural constraint: the uniform allocation of training interventions to heterogeneous farmer populations. Farmers differ systematically in baseline competencies, risk perception, resource access, and learning readiness, yet most training programs continue to apply standardized curricula across all participants. This mismatch results in inefficient resource use and suboptimal learning outcomes, highlighting the need for more adaptive and intelligent training allocation mechanisms.

Recent progress in machine learning (ML) enables the modeling of complex, non-linear relationships within agricultural datasets, offering new opportunities to predict heterogeneous outcomes across individuals. However, predictive capability alone is insufficient for real-world deployment. Decision-making in agricultural extension requires transparency, interpretability, and actionable logic that can be operationalized by practitioners. This has led to increasing interest in explainable artificial intelligence (XAI), particularly SHapley Additive exPlanations (SHAP), which provides theoretically grounded feature attribution for tree-based models. Despite these advances, a critical gap remains: the lack of integrated AI frameworks that translate predictive and explainable insights into structured, deployable decision support systems for agricultural training design.

This study addresses this gap by proposing an explainable AI-based Decision Support System (DSS) for precision allocation of agricultural skill development programs. The framework integrates K-Means clustering for farmer segmentation, Random Forest modeling for skill outcome prediction, and SHAP analysis for interpretable attribution, organized within a unified Segment–Attribute–Recommend pipeline. Unlike prior studies that treat explainability as a post-hoc interpretive tool, this study operationalizes SHAP outputs as the foundation for generating explicit decision rules, enabling the transition from predictive analytics to actionable decision intelligence.

The study is guided by three research questions: (1) Can unsupervised learning methods identify structurally meaningful farmer segments that are relevant for training design? (2) Can machine learning models provide sufficient predictive performance for individual-level skill improvement to support decision-making? and (3) how can SHAP-based explanations be transformed into interpretable and implementable decision rules for agricultural extension systems?

The contributions of this study are threefold. First, it advances AI in agriculture by demonstrating how explainable machine learning can be integrated into an operational decision framework for human capital development, expanding the scope of AI applications beyond physical and production-oriented systems. Second, it introduces cluster-level SHAP analysis as a mechanism for identifying heterogeneous drivers of skill development across farmer groups, enabling more granular and context-specific intervention design. Third, it proposes a deployable DSS architecture that converts model outputs into transparent and auditable decision rules, facilitating real-world implementation by extension officers without requiring direct interaction with underlying AI models.

Importantly, this study does not aim to establish causal relationships between training interventions and outcomes. The dataset is observational and non-experimental in nature; therefore, all findings are interpreted as associative. The purpose of the framework is not to infer causal impact, but to support data-driven decision-making under real-world constraints where experimental control is limited.

The remainder of the paper is structured as follows. Section 2 reviews relevant literature on farmer heterogeneity, machine learning in agriculture, and explainable AI. Section 3 describes the dataset and methodological framework. Section 4 presents the empirical results across segmentation, prediction, and explanation components. Section 5 discusses the implications for AI-driven agricultural decision systems. Section 6 concludes and outlines directions for future research.

\end{abstract}

%%Graphical abstract
% \begin{graphicalabstract}
% \includegraphics{grabs}
% \end{graphicalabstract}

%%Research highlights
% \begin{highlights}
% \item Research highlight 1
% \item Research highlight 2
% \end{highlights}

%% Keywords
\begin{keyword}
Explainable AI \sep Decision Support System \sep SHAP \sep Machine Learning \sep Smart Agriculture
\end{keyword}

\end{frontmatter}

%% Add \usepackage{lineno} before \begin{document} and uncomment 
%% following line to enable line numbers
%% \linenumbers

%% main text
%%

%% Use \section commands to start a section
\section{Introduction}
The rapid advancement of artificial intelligence (AI) has transformed agricultural systems from input-driven production models into data-driven decision ecosystems. While smart agriculture technologies have primarily focused on optimizing physical processes such as crop monitoring, irrigation, and yield prediction, comparatively less attention has been devoted to the development of AI-driven systems that support human capital formation---a critical determinant of long-term agricultural productivity and resilience.

Agricultural skill development programs represent one of the most widely implemented instruments for improving farmer performance. However, their effectiveness remains constrained by the uniform allocation of training interventions across heterogeneous farmer populations. Farmers differ systematically in baseline competencies, risk perception, resource access, and learning readiness, yet standardized curricula are commonly applied across all participants. This mismatch results in inefficient resource allocation and suboptimal learning outcomes, highlighting the need for more adaptive and intelligent allocation mechanisms \cite{anderson2004,feder1985}.

Recent advances in machine learning (ML) \cite{kamilaris2018,liakos2018,yang2025} enable the modeling of complex and non-linear relationships within agricultural datasets, offering new opportunities to predict heterogeneous outcomes across individuals. However, predictive capability alone is insufficient for real-world deployment. Decision-making in agricultural extension requires transparency, interpretability, and actionable logic that can be operationalized by practitioners. This has led to increasing interest in explainable artificial intelligence (XAI)\cite{lundberg2017,arrieta2020,ryo2022}, particularly SHapley Additive exPlanations (SHAP), which provides theoretically grounded feature attribution for tree-based models \cite{lundberg2017}. Despite these advances, a critical gap remains: the lack of integrated AI frameworks that translate predictive and explainable insights into structured and deployable decision support systems for agricultural training design.

Artificial intelligence has been increasingly applied to automate decision-making processes in agriculture, including quality inspection, prediction, and classification tasks \cite{yang2025,wonggasem2024}.

Deep learning models have demonstrated superior performance compared to traditional approaches in agricultural image-based applications \cite{wonggasem2024}.

This study addresses this gap by proposing an explainable AI-based Decision Support System (DSS) for precision allocation of agricultural skill development programs. The framework integrates K-Means clustering for farmer segmentation, Random Forest modeling for skill outcome prediction, and SHAP analysis for interpretable attribution, organized within a unified Segment--Attribute--Recommend pipeline. Unlike prior studies that treat explainability as a post-hoc interpretive tool, this study operationalizes SHAP outputs as the foundation for generating explicit decision rules, enabling the transition from predictive analytics to actionable decision intelligence.

The study is guided by three research questions: (1) Can unsupervised learning methods identify structurally meaningful farmer segments for training design? (2) Can machine learning models provide sufficient predictive performance to support individual-level decision-making in skill development? and (3) how can SHAP-based explanations be transformed into interpretable and implementable decision rules for agricultural extension systems?

This study makes three primary contributions to the field of artificial intelligence in agriculture. First, it extends the role of AI beyond predictive analytics toward actionable decision-making by proposing an integrated explainable AI-based decision support framework for human capital development. Second, it advances explainable AI by introducing a cluster-level SHAP approach that links model interpretability with structured decision logic. Third, it proposes a deployable Segment--Attribute--Recommend architecture that transforms model outputs into transparent and auditable decision rules, enabling practical implementation by domain practitioners.

Importantly, this study does not aim to establish causal relationships between training interventions and outcomes. The dataset is observational and non-experimental in nature; therefore, all findings are interpreted as associative. The framework is intended to support data-driven decision-making under real-world constraints rather than causal inference.

The remainder of the paper is structured as follows. Section~2 reviews the relevant literature. Section~3 presents the data and methodology. Section~4 reports the empirical results. Section~5 discusses the implications for AI-driven agricultural systems. Section~6 concludes the study.
\label{sec1}
%% Labels are used to cross-reference an item using \ref command.

Section text. See Subsection \ref{subsec1}.

%% Use \subsection commands to start a subsection.
\subsection{Example Subsection}
\label{subsec1}

Subsection text.

%% Use \subsubsection, \paragraph, \subparagraph commands to 
%% start 3rd, 4th and 5th level sections.
%% Refer following link for more details.
%% https://en.wikibooks.org/wiki/LaTeX/Document_Structure#Sectioning_commands

\subsubsection{Mathematics}
%% Inline mathematics is tagged between $ symbols.
This is an example for the symbol $\alpha$ tagged as inline mathematics.

%% Displayed equations can be tagged using various environments. 
%% Single line equations can be tagged using the equation environment.
\begin{equation}
f(x) = (x+a)(x+b)
\end{equation}

%% Unnumbered equations are tagged using starred versions of the environment.
%% amsmath package needs to be loaded for the starred version of equation environment.
\begin{equation*}
f(x) = (x+a)(x+b)
\end{equation*}

%% align or eqnarray environments can be used for multi line equations.
%% & is used to mark alignment points in equations.
%% \\ is used to end a row in a multiline equation.
\begin{align}
 f(x) &= (x+a)(x+b) \\
      &= x^2 + (a+b)x + ab
\end{align}

\begin{eqnarray}
 f(x) &=& (x+a)(x+b) \nonumber\\ %% If equation numbering is not needed for a row use \nonumber.
      &=& x^2 + (a+b)x + ab
\end{eqnarray}

%% Unnumbered versions of align and eqnarray
\begin{align*}
 f(x) &= (x+a)(x+b) \\
      &= x^2 + (a+b)x + ab
\end{align*}

\begin{eqnarray*}
 f(x)&=& (x+a)(x+b) \\
     &=& x^2 + (a+b)x + ab
\end{eqnarray*}

%% Refer following link for more details.
%% https://en.wikibooks.org/wiki/LaTeX/Mathematics
%% https://en.wikibooks.org/wiki/LaTeX/Advanced_Mathematics

%% Use a table environment to create tables.
%% Refer following link for more details.
%% https://en.wikibooks.org/wiki/LaTeX/Tables
\begin{table}[t]%% placement specifier
%% Use tabular environment to tag the tabular data.
%% https://en.wikibooks.org/wiki/LaTeX/Tables#The_tabular_environment
\centering%% For centre alignment of tabular.
\begin{tabular}{l c r}%% Table column specifiers
%% Tabular cells are separated by &
  1 & 2 & 3 \\ %% A tabular row ends with \\
  4 & 5 & 6 \\
  7 & 8 & 9 \\
\end{tabular}
%% Use \caption command for table caption and label.
\caption{Table Caption}\label{fig1}
\end{table}


%% Use figure environment to create figures
%% Refer following link for more details.
%% https://en.wikibooks.org/wiki/LaTeX/Floats,_Figures_and_Captions
\begin{figure}[t]%% placement specifier
%% Use \includegraphics command to insert graphic files. Place graphics files in 
%% working directory.
\centering%% For centre alignment of image.
\includegraphics{example-image-a}
%% Use \caption command for figure caption and label.
\caption{Figure Caption}\label{fig1}
%% https://en.wikibooks.org/wiki/LaTeX/Importing_Graphics#Importing_external_graphics
\end{figure}


%% The Appendices part is started with the command \appendix;
%% appendix sections are then done as normal sections
\appendix
\section{Example Appendix Section}
\label{app1}

Appendix text.

%% For citations use: 
%%       \cite{<label>} ==> [1]

%%
Example citation, See \cite{lamport94}.

%% If you have bib database file and want bibtex to generate the
%% bibitems, please use
%%
%%  \bibliographystyle{elsarticle-num} 
%%  \bibliography{<your bibdatabase>}

%% else use the following coding to input the bibitems directly in the
%% TeX file.

%% Refer following link for more details about bibliography and citations.
%% https://en.wikibooks.org/wiki/LaTeX/Bibliography_Management

\begin{thebibliography}{00}

%% For numbered reference style
%% \bibitem{label}
%% Text of bibliographic item

\bibitem{lamport94}
  Leslie Lamport,
  \textit{\LaTeX: a document preparation system},
  Addison Wesley, Massachusetts,
  2nd edition,
  1994.

\end{thebibliography}
\end{document}

\endinput
%%
%% End of file `elsarticle-template-num.tex'.
